
\section{Bounded localisation and the corrected absorption theorem}\label{app:bounded-localisation}

We include the bounded quotient argument to make the hypotheses used in the main text explicit.

\begin{theorem}[Bounded localisation equivalence]\label{thm:appendix-bounded-localisation}
Assume \(\lambda:R\to S\) is a homological ring epimorphism satisfying \textup{(B1)--(B3)}.  Then extension of scalars induces an exact equivalence
\[
 \Db(R\text{-}\modcat)/\Kb(\lambda)
 \xrightarrow{\sim}
 \Db(S\text{-}\modcat).
\]
\end{theorem}

\begin{proof}
The unbounded adjunction is
\[
 \Phi=S\Ltensor_R-:D(R)\rightleftarrows D(S):U.
\]
By \Cref{lem:bounded-adjunction}, this adjunction restricts to the bounded finitely generated subcategories.  Homologicality makes \(U\) fully faithful and the counit \(\Phi U\to\id\) an isomorphism.

Let \(q:\Db(R\text{-}\modcat)\to\Db(R\text{-}\modcat)/\Kb(\lambda)\).  For every bounded \(X\), complete the unit to
\[
 K_X\longrightarrow X\longrightarrow U\Phi X\longrightarrow K_X[1].
\]
Applying \(\Phi\) and using the counit shows \(\Phi K_X=0\); boundedness of the other two terms puts \(K_X\) in \(\Kb(\lambda)\).  Therefore \(qX\cong qU\Phi X\).  The counit gives the inverse natural isomorphism on the \(S\)-side.  Thus the induced quotient functor is an equivalence.
\end{proof}

For a finitely generated module \(M\), write
\[
 \Dsg(R,M)=\Db(R\text{-}\modcat)/\thick_{\Db(R)}(M).
\]
This definition avoids any implicit claim that a homotopy category of complexes with terms in \(\add M\) embeds fully in the derived category.

\begin{theorem}[Correct bounded absorption theorem]\label{thm:correct-absorption}
Assume the hypotheses of \Cref{thm:appendix-bounded-localisation}.  Let \(M\) be a finitely generated \(R\)-module and assume
\[
 \Tor_i^R(S,X)=0
 \qquad(i>0,\ X\in\add M).
\]
If
\[
 \Kb(\lambda)\subseteq\thick_{\Db(R)}(M),
\]
then extension of scalars induces an exact equivalence
\[
 \Dsg(R,M)
 \xrightarrow{\sim}
 \Dsg(S,S\otimes_RM).
\]
\end{theorem}

\begin{proof}
The Tor vanishing identifies \(S\Ltensor_RM\) with the ordinary module \(S\otimes_RM\).  Exactness implies
\[
 \Phi\bigl(\thick(M)\bigr)
 \subseteq
 \thick(S\otimes_RM).
\]
Conversely, the latter thick subcategory is generated by the image of \(M\), so it is exactly the thick closure of the image of \(\thick(M)\) under the bounded localisation equivalence.

Since \(\Kb(\lambda)\subseteq\thick(M)\), the third isomorphism theorem gives
\[
 \Db(R)/\thick(M)
 \simeq
 \bigl(\Db(R)/\Kb(\lambda)\bigr)
 \Big/
 \bigl(\thick(M)/\Kb(\lambda)\bigr).
\]
Use \Cref{thm:appendix-bounded-localisation} and identify the denominator with \(\thick(S\otimes_RM)\).
\end{proof}

\begin{corollary}[Defect absorption]\label{cor:defect-absorption}
Under the hypotheses above, suppose
\[
 \Kc_\lambda\subseteq\thick(M)
\]
and the induced functor
\[
 \DeltaB(\lambda)
 \longrightarrow
 \Db(R\text{-}\modcat)/\thick(M)
\]
is zero.  Then \(\Kb(\lambda)\subseteq\thick(M)\), so
\Cref{thm:correct-absorption} applies.
\end{corollary}

\begin{proof}
The quotient functor on \(\Kb(\lambda)\) factors through
\(\Kb(\lambda)/\Kc_\lambda=\DeltaB(\lambda)\).  Its induced functor is zero by assumption, hence every bounded kernel object maps to zero in the quotient by \(\thick(M)\).
\end{proof}
